\section{The common moment information in the Fourier bounds}
\label{ext:fourier-moments}
The finite Fourier estimates use two kinds of information. Bounds valid for
every standardized law cover the larger moment ranges. Smaller moment ranges
also use the support bound already obtained from contact. The distinction is
essential: a support hypothesis cannot be inferred from a successful numerical
integral.

Let $X$ be centered with $\E X^2=1$, $\rho=\E|X|^3<\infty$,
$\gamma=\E X^3$, and $f(u)=\E e^{iuX}$. In this section put
\[
 \begin{gathered}
 h=\E|X|,\quad m=\E X_+=\E X_-=h/2,\quad
 A=\E X_+^2,\quad B=\E X_-^2,\\
 \zeta=2AB/m,\qquad b(r)=\frac{\sqrt{r^2+8}-r}{2}.
 \end{gathered}
\]
Here $X_-=\max\{-X,0\}$, so $A+B=1$ and all of $h,m,A,B$ are
positive. The letter $b$ is used for this root function to avoid confusion
with the signed second moment $\beta=\E[X|X|]$.

\begin{lemma}[Anchored characteristic function]\label{ext:fourier-anchor}
For $0\le u\le\pi$,
\[
 \cos u\le\Re f(u)\le\cos u+(\rho-1)u^3/6,
 \qquad
 |\Im f(u)|^2\le\frac{u^6}{36}(\rho-1)(\rho+5/3).
\]
\end{lemma}
\begin{proof}
Write $\operatorname{sinc}t=\int_0^1\cos(st)\,ds$. For real $a,b$,
\[
 |\operatorname{sinc}a-\operatorname{sinc}b|
 \le |a-b|/2,\qquad
 |\operatorname{sinc}a-\operatorname{sinc}b|
 \le |a^2-b^2|/6.
\]
The first follows from the Lipschitz bound for cosine. For the second,
$v\mapsto\cos\sqrt v$ has derivative of absolute value at most $1/2$
on $v\ge0$, including its limit at zero; integrate with respect to $s$.
For $r\ge0$, integration from $1$ to $r$ of the first estimate gives
\[
 \left|\cos(ur)-\cos u+\frac{u\sin u}{2}(r^2-1)\right|
 \le \frac{u^3}{12}(r-1)^2(2r+1).
\]
Since $\E(|X|-1)^2(2|X|+1)=2(\rho-1)$, averaging proves the
upper real bound. For the lower bound, when $0<u\le\pi$ set
$a=\sin u/(2u)$ and $J(v)=\cos v+av^2$. The function
$\operatorname{sinc}v$ decreases on $[0,\pi]$, so $J$ decreases up to
$u$ and increases from $u$ to $\pi$. For $v\ge\pi$,
$J(v)\ge-1+a\pi^2=J(\pi)\ge J(u)$. Evenness gives the same conclusion
for negative $v$. Average $J(uX)\ge J(u)$ and use $\E X^2=1$.

The second sinc estimate, together with $\E X=0$, gives
$|\Im f(u)|\le u^3\E[|X|\,||X|^2-1|]/6$.
Put $W(r)=(r-1)^2(2r+1)$. Cauchy's inequality applies to the factorization
$r|r^2-1|=\sqrt{W(r)}\,r(r+1)/\sqrt{2r+1}$.
The exact identity
\[
 r^3+\frac{29r^2+1}{18}
 -\frac{2r^2(r+1)^2}{2r+1}
 =\frac{(r-1)^2(4r+1)}{18(2r+1)}\ge0
\]
bounds the second squared factor in expectation by $(\rho+5/3)/2$;
the first has expectation $2(\rho-1)$. All expectations are finite under
the stated third-moment hypothesis. The value at $u=0$ is immediate.
\end{proof}

\begin{lemma}[One common moment surrogate]\label{ext:fourier-surrogate}
For an independent copy $X'$, put
$\eta=\E|X-X'|^3/2-\rho$. Then
\[
 \eta\le h\le1,\qquad \zeta\ge h,\qquad
 \rho+\zeta\ge2/h,\qquad \zeta\ge b(\rho),
 \qquad \gamma^2+\zeta^2\le\rho^2.
\]
For every real $u$ one also has
$36|\Im f(u)|^2\le u^6(\rho^2-\zeta^2)$.
If $|X|\le R$ almost surely, then additionally
\[
 \eta\le b(\rho)+\frac{(R-\sqrt{\rho^2+8})_+^2}{4R}.
\]
In particular $R\le\sqrt{\rho^2+8}$ implies $\eta\le b(\rho)$.
\end{lemma}
\begin{proof}
The scalar inequality
$|x-y|^3\le(x-y)^2(|x|+|y|)$, independence and centering give
$\E|X-X'|^3\le2\rho+2h$; Cauchy gives $h\le1$.
Signwise Cauchy gives
\[
 A^2\le m\E X_+^3,\quad B^2\le m\E X_-^3,
 \quad \frac{m^2}{A}+\frac{m^2}{B}\le\Pp(X\ne0)\le1.
\]
Thus $AB\ge m^2$, $\zeta\ge2m=h$, and
$\rho+\zeta\ge(A+B)^2/m=2/h$. Multiplication yields
$\zeta(\rho+\zeta)\ge2$, proving $\zeta\ge b(\rho)$.

For completeness, the sine-circle bound uses an exact mixture, not an
extreme-point assertion about convolution. Let $\mu_+$ and $\mu_-$ be the
restrictions of the law of $|X|$ to its positive and negative sign classes.
On $a,b>0$ define
\[
 d\nu(a,b)=\frac{a+b}{m}\,d\mu_+(a)d\mu_-(b).
\]
Integrating the centered law with atoms $a,-b$ and probabilities
$b/(a+b),a/(a+b)$ against $\nu$ recovers the nonzero part of the original
law; the zero atom is added separately. For this two-point law the second
sinc estimate implies
\[
 \frac{6|\Im f_{a,b}(u)|}{|u|^3}
 \le\frac{ab|a^2-b^2|}{a+b}.
\]
Together with $e_{a,b}=2a^2b^2/(a+b)$, the right side is one coordinate
of a vector of length $ab(a^2+b^2)/(a+b)$. Minkowski's inequality in
$\R^2$ therefore bounds the vector
$(\zeta,6\Im f(u)/u^3)$ by $\rho$. The third-moment assumption
justifies all integrals. Letting $u\to0$ proves the assertion for $\gamma$:
Taylor's remainder is dominated by $|X|^3/6$, so dominated convergence
applies after using $\E X=0$.

For the final assertion the exact same-sign deficit is
\[
 h-\eta=\E\!\left[
 \min(|X|,|X'|)(X-X')^2\ind_{\{XX'>0\}}\right].
\]
On a common sign class, $\min(a,b)\ge ab/R$. Expansion of the two
independent class integrals gives
\[
 h-\eta\ge\frac{h\rho-2(A^2+B^2)}{R}
 =\frac{h(\rho+\zeta)-2}{R}.
\]
If $h\le b(\rho)$, the already proved $\eta\le h$ suffices. Otherwise,
write $b=b(\rho)$, use $\zeta\ge h$ and $b(\rho+b)=2$ to obtain
\[
 \eta-b\le\frac{(h-b)(R-\rho-h-b)}{R}\le0,
\]
when $R\le\rho+2b=\sqrt{\rho^2+8}$. For general $R$, the quadratic
$y(R-\sqrt{\rho^2+8}-y)$ on $y\ge0$ is at most
$(R-\sqrt{\rho^2+8})_+^2/4$. This proves the stated remainder too.
In particular radius three suffices for every $\rho\ge1$.
\end{proof}

For a universal calculation put $H_* =\max\{h,b(\rho)\}$.
Then $b(\rho)\le H_*\le1$, $\eta\le H_*$, $\zeta\ge H_*$ and
$\rho+\zeta\ge2/H_*$. On a closed cell
$\rho\in[\ell,r]$, $H_*\in[a,b]$, it is consequently safe to use
\[
 \eta\le b,\qquad
 \zeta\ge\max\{a,2/b-r,b(r)\}.
\]
Every cell uses both bounds from the same marginal law. A lower bound on
$H_*$ is not a lower bound on $h$; no later sign-probability argument may
substitute it for one. If support $|X|\le R\le\sqrt{\ell^2+8}$ has independently been acquired,
the last part of Lemma~\ref{ext:fourier-surrogate} instead permits
$\eta\le b(\ell)$ and $\zeta\ge b(r)$. For larger $R$ its displayed
nonnegative remainder must be retained.

\subsection{Keeping the cubic compensation in the high frequencies}
\label{ext:compensated-high}
Suppose now that $|X|\le R$ and $\rho\le r$. We may use
\begin{align*}
 M_4&=\min\left\{Rr,1+\frac{(R+1)^2}{R+1/2}(r-1)\right\},\\
 M_5&=\min\left\{RM_4,
 1+\frac{R^3+2R^2+3R+3/2}{R+1/2}(r-1)\right\}
\end{align*}
as upper bounds on $\E X^4$ and $\E|X|^5$ of the same law. To verify the
less immediate branches, set $y=|X|$ and expand
\[
 \rho-1=\E(y-1)^2(y+1/2),\quad
 \E X^4-1=\E(y-1)^2(y+1)^2,
\]
\[
 \E|X|^5-1=\E(y-1)^2(y^3+2y^2+3y+3/2).
\]
The ratios of the last two factors to $y+1/2$ increase on $y\ge0$:
their derivatives are respectively
$y(y+1)/(y+1/2)^2$ and $y(4y^2+7y+4)/(2(y+1/2)^2)$.
The other branches follow from $y\le R$. These moment bounds require
no bound on the product-gap defect.

Put $G=\min\{(r-1)(r+5/3),r^2-Z^2\}$, where $0<Z\le b(r)$.
Assume that a proved nonnegative envelope $m(u)$ bounds $|f(u)|$ on
the interval being integrated. Set
\[
 g=e^{-u^2/2},\quad v=-i\gamma u^3/6,\quad
 b_u=\sqrt G\,u^3/6,\quad W=g\sqrt{1+b_u^2},
\]
and, for $0\le u\le\pi$,
\[
 \begin{gathered}
 A_u=\max\{|\cos u-g|,
 |\cos u+(r-1)u^3/6-g|\},\\
 B_u=\min\{M_4u^4/24, M_5u^5/120\}+b_u(1-g).
 \end{gathered}
\]
Taylor's theorem gives
\[
 |\Im f(u)+\gamma u^3/6|\le\min\{M_4u^4/24,M_5u^5/120\}.
\]
The anchored real estimate therefore gives
$|f-g(1+v)|\le\sqrt{A_u^2+B_u^2}$. Telescoping the $n$th powers and
then applying the integral Taylor remainder to $(1+sv)^n$, $0\le s\le1$,
proves, for every integer $n\ge2$,
\begin{align*}
 |f(u)^n-g^n(1+nv)|
 &\le\sqrt{A_u^2+B_u^2}\sum_{j=0}^{n-1}m(u)^jW^{n-1-j}
 +\binom n2 b_u^2g^n(1+b_u^2)^{(n-2)/2},\\
 |f(u)^n-g^n(1+nv)|
 &\le m(u)^n+g^n\sqrt{1+n^2b_u^2}.
\end{align*}
One may take their pointwise minimum. These formulas use no quotient of
$m(u)-W$, so coincident and zero values cause no exceptional cases.

On a closed frequency cell $[a,b]\subset[1/2,T]$, a rigorous upper
bound for this minimum is multiplied by the exact weight
$\int_a^b u^{-3}\,du=(a^{-2}-b^{-2})/2$. The infinite remainder is
bounded without a periodicity assumption: for $T>0$,
\[
 \int_T^\infty\frac{|f(u)^n-g(u)^n(1+nv(u))|}{u^3}\,du
 \le\frac1{2T^2}
 +\frac{e^{-nT^2/2}}{nT}
 \left(T^{-3}+\frac{n\sqrt G}{6}\right).
\]
Indeed $|f|\le1$, $\sqrt{1+x^2}\le1+|x|$, $u^{-3}\le T^{-3}$
on the normal part, and
$\int_T^\infty e^{-nu^2/2}\,du\le e^{-nT^2/2}/(nT)$.
This explains both the compensation and the complete tail in the native
high-frequency certificates.

To preserve cancellation in the real residual on a closed frequency cell,
put $d_0(u)=\cos u-e^{-u^2/2}$ and
$d_1(u)=d_0(u)+(r-1)u^3/6$. Their derivatives are
\[
 d_0'(u)=-\sin u+ue^{-u^2/2},\qquad
 d_1'(u)=d_0'(u)+(r-1)u^2/2.
\]
If $D_j$ encloses $d_j'$ throughout a cell of width $w$ and midpoint
$u_0$, the fundamental theorem of calculus gives
$d_j(u)\in d_j(u_0)+[-w/2,w/2]D_j$.
Intersecting this with a direct interval evaluation remains valid.
This avoids losing the cancellation in $\cos u-e^{-u^2/2}$ by bounding
its two terms separately.

\subsection{The compatible modulus profile}
\label{ext:cosine-profile}
The scalar certificate uses $\kappa=62/625$ and proves
\begin{equation}\label{ext:scalar-cosine}
 \cos x\le1-x^2/2+\kappa|x|^3\qquad(x\in\R).
\end{equation}
Here is the entire reduction to its finite check. For $x\ge625/124$ the
right side is at least one. On the remaining interval, the degree-sixteen
Taylor polynomial $P_{16}$ bounds cosine above: the sixteenth derivative
of $\cos-P_{16}$ is $\cos-1\le0$, and its lower derivatives vanish at
zero. The degree-thirteen polynomial
\[
 Q(x)=\kappa+\sum_{j=2}^{8}\frac{(-1)^{j+1}x^{2j-3}}{(2j)!}
\]
satisfies $x^3Q(x)=1-x^2/2+\kappa x^3-P_{16}(x)$.
On the eighty cells $[i/16,(i+1)/16]$, $0\le i<80$, and the last
cell $[5,625/124]$, the exact verifier checks the inverse Bernstein
identity and all fourteen coefficients. The smallest coefficient is
$2479189/65585520000>0$. Lemma~\ref{ext:bernstein} proves positivity
on every closed cell; evenness completes \eqref{ext:scalar-cosine}.
The native checker rejects twenty-one malformed objects and also checks
324 exact evaluations of the power/Bernstein identity.

Put $a=1/2-4\kappa=129/1250$, $C=1-\cos4$ and $A_0=16a/C$.
The profile
\[
 q(z)=\begin{cases}
 1/2-\kappa z,&0\le z\le4,\\
 A_0(1-\cos z)/z^2,&4\le z\le2\pi,\\
 0,&z\ge2\pi
 \end{cases}
\]
is nonnegative, convex and nonincreasing on $[0,\infty)$, and satisfies
$\cos z\le1-z^2q(z)$. Indeed \eqref{ext:scalar-cosine} gives
$0<A_0\le1$. The pieces agree at their endpoints. On the middle interval
the first derivative is nonpositive and the second has the sign of
\[
 B_0(z)=6+(z^2-6)\cos z-4z\sin z.
\]
On $[4,3\pi/2]$, $B_0'(z)=-2z\cos z+(2-z^2)\sin z\ge0$ and
$B_0(4)\ge4$: indeed $7\pi/6\le4\le3\pi/2$ implies
$\sin4\le-1/2$, and $\cos4\ge-1$, so
$6+10\cos4-16\sin4\ge4$. On $[3\pi/2,2\pi]$, $B_0(z)\ge6$.
The derivative at $2\pi$ is zero. The derivative jump at $4$ has the
correct sign because $258\sin4+119(1-\cos4)>0$. Its exact lower bound is
obtained by replacing sine by its degree-fifteen Taylor lower bound and
cosine by $P_{16}$. The same scalar checker verifies this rational
inequality and the needed $25/8<\pi<22/7$ by Machin's formula.

To apply the profile, weight the distribution of $X-X'$ by
$(X-X')^2/2$. This is a probability measure whose first absolute moment
is $\rho+\eta$. Jensen's inequality gives, for $u\ge0$,
\[
 |f(u)|^2=\E\cos(u(X-X'))
 \le1-2u^2q(u(\rho+\eta)).
\]
If $\rho+\eta\le D$, monotonicity permits $D$ on the right. Thus the
linear part gives $1-u^2+2\kappa Du^3$; when $4\le Du\le2\pi$ the
curved part gives $1-2A_0(1-\cos(Du))/D^2$.
Take the minimum with one and with the coordinate-disk estimate from
Lemmas~\ref{ext:fourier-anchor}--\ref{ext:fourier-surrogate}. The upper
square of a real interval is the larger of the squares of its two
endpoints, even when that interval crosses zero or lies below it.

\subsection{From characteristic functions to the CDF}
\label{ext:finite-fourier-transfer}
We use Prawitz's smoothing inequality in the form recorded in
I.~Shevtsova, \emph{Annales Mathematicae et Informaticae} 39 (2012),
Lemma~4.1, p.~263. For $n\ge2$, $c>0$, $0<\tau\le1$, put
\[
 K_2(t)=(1-t)+i\{(1-t)\cot(\pi t)+1/\pi\},\qquad 0<t<1.
\]
With $T=c\sqrt n$, that inequality gives
\begin{align*}
 \Delta_n\le{}&\int_0^\tau |K_2(t)|
   |f(ct)^n-e^{-nc^2t^2/2}|\,dt
 +\int_\tau^1|K_2(t)|\,|f(ct)|^n\,dt\\
 &+\int_0^\tau|K_2(t)-i/(\pi t)|e^{-nc^2t^2/2}\,dt
 +\frac1\pi\int_\tau^\infty e^{-nc^2t^2/2}\frac{dt}{t}.
\end{align*}
Our standardized iid laws satisfy its independence, centering and positive
total-variance hypotheses. The cited convention uses a strict CDF.
Continuity of $\Phi$ makes the supremum identical for closed CDFs: take
one-sided limits at each evaluation point. Thus atoms require no smoothing
of the original law.

For $g(u)=e^{-u^2/2}$ and envelopes $m\ge|f|$, $E_0\ge|f-g|$,
the first integrand uses
$|f^n-g^n|\le\min\{E_0\sum_{j=0}^{n-1}m^jg^{n-1-j},m^n+g^n\}$.
Every literal integer $n$ is evaluated separately. The original $20$--$24$ integral verifier
covers $[\epsilon,\tau]$ and $[\tau,1-\epsilon]$ by closed cells,
where $\epsilon=1/1024$; its numerical parameter domain is
$0<c<\pi$ and $\epsilon<\tau\le3/10<1-\epsilon$.
In particular $c\tau<1$ licenses the stable low-frequency Taylor bounds.
On a cell of width $w$ and midpoint $t_0$,
it uses either $w\sup F$ or
$w\{F(t_0)+w^2\max(0,\sup F'')/24\}$ for a fixed smooth majorant
$F$. The latter follows by Taylor expansion and integration; a failed
derivative-domain guard triggers the former bound on the entire cell.
Switching branches at a sampled midpoint would not justify this rule.

For the omitted endpoint pieces, the respective upper bounds are
\begin{align*}
 L_\epsilon&=n\left[
 \frac{rc^3}{6}\left(\frac{\epsilon^4}{4}
       +\frac{\epsilon^3}{3\pi}\right)
 +\frac{c^4}{8}\left(\frac{\epsilon^5}{5}
       +\frac{\epsilon^4}{4\pi}\right)\right],\\
 N_\epsilon&=\epsilon+\pi\epsilon^2/2,\qquad
 H_\epsilon=\epsilon^2/2+\pi\epsilon^3/3.
\end{align*}
They follow from $|K_2(t)|\le1+1/(\pi t)$ near zero, cubic Taylor
error for $f$, $e^{-u^2/2}-1+u^2/2\le u^4/8$, and
\[
 0\le1/x-\cot x\le\frac{x}{3(1-x^2/6)}\le x\quad(0<x\le1).
\]
For this last estimate use
$\sin x-x\cos x=\int_0^x t\sin t\,dt\le x^3/3$ and
$\sin x\ge x-x^3/6$. Apply the same identity with $1-t$ near one.
Finally the infinite Gaussian term is at most
$e^{-nc^2\tau^2/2}/(\pi nc^2\tau^2)$, by
$\int_a^\infty e^{-v}v^{-1}\,dv\le e^{-a}/a$.
For a strip $\rho\in[\ell,r]$, multiply the full four-term bound by
$\sqrt n/\ell$, with outward enclosures. A complete closed strip cover
and a strict rational comparison then establish the claimed CDF bound.


\subsection{Sharper scalar bounds for the same integrals}
\label{ext:finite-fourier-scalars}
For $a>0$ write $E_1(a)=\int_a^\infty e^{-v}v^{-1}\,dv$.
Integration by parts, followed by Jensen's inequality for the convex
function $t\mapsto(a+t)^{-2}$ and an exponential random variable of mean
one, gives
\[
 E_1(a)=\frac{e^{-a}}a-e^{-a}\int_0^\infty
                \frac{e^{-t}}{(a+t)^2}\,dt
 \le e^{-a}\left(\frac1a-\frac1{(a+1)^2}\right).
\]
The last Gaussian term in Prawitz's inequality is $E_1(a)/(2\pi)$,
where $a=nc^2\tau^2/2$. This improvement has no additional moment or
support hypothesis.

Some small-index certificates require $\tau>3/10$. Their low-frequency
range is still controlled explicitly. Put $U=6/5$. On $0\le x\le U$,
the alternating sine bounds and integration give
\[
 \sin x\ge x\left(1-x^2/6+x^4/120-x^6/5040\right),\qquad
 0\le\sin x-x\cos x\le x^3/3-x^5/30+x^7/840.
\]
The denominator below is at least $1-U^2/6-U^6/5040>0$, and the
numerator factor is at least $1/3-U^2/30>0$. Consequently
\[
 0\le1/x-\cot x\le L(x):=
 \frac{x(1/3-x^2/30+x^4/840)}
      {1-x^2/6+x^4/120-x^6/5040},
 \qquad L(0)=0.
\]
It follows that
\[
 |K_2(t)-i/(\pi t)|
 \le(1-t)\sqrt{1+L(\pi t)^2}
 \quad(0<t\le19/50).
\]
At the high endpoint put $s=1-t$. The identity
$K_2(t)=s+i\,s(1/(\pi s)-\cot(\pi s))$ gives
$|K_2(t)|\le s\sqrt{1+L(\pi s)^2}$ for $t\ge3/4$.
The rational upper bound $\pi<22/7$ implies
$\pi(19/50)<6/5$ and $\pi/4<6/5$.
The numerical quadrature must therefore declare
$\epsilon<\tau\le19/50<1-\epsilon$, $c<\pi$, and $c\tau\le6/5$.
It differentiates the displayed smooth rational majorant on whole cells;
it does not differentiate an interval-selected minimum.

The cancellation series on this larger range also have explicit tails.
Truncate the even series for cosine and $e^{-u^2/2}$ at degree twenty.
For their difference, the error in each $t$ derivative of order
$j=0,1,2$, with $u=ct\in[0,U]$, is at most
\[
 \frac{2(22)_j c^j u_{\max}^{22-j}}{2^{11}11!}.
\]
Indeed, for terms of Gaussian index $k\ge11$, the ratios of consecutive
absolute derivative terms are bounded respectively by
\[
 \frac{U^2}{2(k+1)},\qquad \frac{U^2}{2k},\qquad
 \frac{U^2(2k+1)}{2k(2k-1)}.
\]
All decrease with $k$ and are less than $1/8$ at $k=11$.
The cosine ratios are
$U^2/\{(2k+2-j)(2k+1-j)\}<1/8$.
The first cosine tail is $2^{11}11!/22!<1/4$ times the first Gaussian
tail, for every $j$. Summing both geometric tails gives a factor less
than $(1+1/4)/(1-1/8)=10/7<2$.
Uniform absolute convergence justifies these differentiations.
For the removable function $2(\cos u-e^{-u^2/2})/u^3$, divide the series
symbolically, so its value tail is bounded by
$4u_{\max}^{19}/(2^{11}11!)$ and its value at zero is obtained directly
from the series. No numerical division by a nearly zero frequency is
needed.
